% GENERATED FILE. Edit canonical lessons, then run npm run generate:books.
% canonical-source-hash: fnv1a64:41d2f229daaf4737
% canonical-lessons: ML-C41-this, ML-C41-that, ML-C41-here, ML-C41-there, ML-C41-who, ML-C41-where, ML-C41-deixis-system

\chapter{Pointing, and Asking}
\label{ch:ml-pointing-and-asking}
\hlchaptermodality{41}

\begin{chapteropening}
\textbf{By the end of this chapter:} \emph{I can point at something in Malayalam and ask who or where it is — and I can say what the six words have in common.}

Say the near word, the far word and the question word in a row, and name the part that changed.
\end{chapteropening}

\section[\texorpdfstring{\ml{ഇത്} (itŭ)}{ഇത് (itŭ)}]{\ml{ഇത്} (itŭ) — this one — the thing near me}

\label{lesson:ML-C41-this}



\begin{quote}
You can already name things. This is how you POINT at them.
\end{quote}

\subsection*{You'll want to know: \ml{ഇത്}}
\textbf{\ml{ഇത്}} — \emph{itŭ} — this one — the thing near me.

Say it, and picture the thing it lands on as you say it. That is the whole word: it does not mean anything on its own, it means whatever the speaker is pointing at.

Its partner, the far one, comes in the next lesson. The two of them differ by one sound at the front — \textbf{i-} for near — and that is not a coincidence. It is the whole system, and the last lesson of this chapter says so.

\begin{cousinweb}
\textbf{\ml{ഇത്}} (\emph{itŭ}) is native Dravidian: the old near-pointing base with an ending meaning \textquotedblleft{}thing\textquotedblright{}. DEDR 410 lists Malayalam \emph{itu} with Tamil \emph{itu}.
\end{cousinweb}

\subsection*{Your turn}
\begin{itemize}
\raggedright
  \item \emph{Say it:} \textquotedblleft{}\ml{ഇത്}\textquotedblright{} three times, picturing something different each time
  \item \emph{Say it:} it once with a word you already know after it
  \item \emph{Recall:} say \emph{cevi}
  \item \emph{Read it:} \textbf{\ml{മൂക്ക്}}
\end{itemize}

\begin{tcolorbox}[breakable,colback=teal!4,colframe=teal!35!black,title={Before you move on}]
What is \emph{itŭ}? (this one — the thing near me.) What sound does it start with, and does that sound mean near, far, or a question? (\textbf{i-} — near.)
\end{tcolorbox}
\section[\texorpdfstring{\ml{അത്} (atŭ)}{അത് (atŭ)}]{\ml{അത്} (atŭ) — that one — the thing over there}

\label{lesson:ML-C41-that}



\begin{quote}
You can already name things. This is how you POINT at them.
\end{quote}

\subsection*{You'll want to know: \ml{അത്}}
\textbf{\ml{അത്}} — \emph{atŭ} — that one — the thing over there.

Say it, and picture the thing it lands on as you say it. That is the whole word: it does not mean anything on its own, it means whatever the speaker is pointing at.

Its partner is \textbf{\ml{ഇത്}} \emph{itŭ}, which you have just met. The two of them differ by one sound at the front — \textbf{a-} for far — and that is not a coincidence. It is the whole system, and the last lesson of this chapter says so.

\begin{cousinweb}
\textbf{\ml{അത്}} (\emph{atŭ}) is native Dravidian: the old far-pointing base with an ending meaning \textquotedblleft{}thing\textquotedblright{}. DEDR 1 lists Malayalam \emph{atu} with Tamil \emph{atu}, Kannada \emph{adu} and Telugu \emph{adi}.
\end{cousinweb}

\subsection*{Your turn}
\begin{itemize}
\raggedright
  \item \emph{Say it:} \textquotedblleft{}\ml{അത്}\textquotedblright{} three times, picturing something different each time
  \item \emph{Say it:} it once with a word you already know after it
  \item \emph{Read it:} \textbf{\ml{വായ്}}
\end{itemize}

\begin{tcolorbox}[breakable,colback=teal!4,colframe=teal!35!black,title={Before you move on}]
What is \emph{atŭ}? (that one — the thing over there.) What sound does it start with, and does that sound mean near, far, or a question? (\textbf{a-} — far.)
\end{tcolorbox}
\section[\texorpdfstring{\ml{ഇവിടെ} (iviṭe)}{ഇവിടെ (iviṭe)}]{\ml{ഇവിടെ} (iviṭe) — here — where I am}

\label{lesson:ML-C41-here}



\begin{quote}
You can already name things. This is how you POINT at them.
\end{quote}

\subsection*{You'll want to know: \ml{ഇവിടെ}}
\textbf{\ml{ഇവിടെ}} — \emph{iviṭe} — here — where I am.

Say it, and picture the thing it lands on as you say it. That is the whole word: it does not mean anything on its own, it means whatever the speaker is pointing at.

Its partner, the far one, comes in the next lesson. The two of them differ by one sound at the front — \textbf{i-} for near — and that is not a coincidence. It is the whole system, and the last lesson of this chapter says so.

\begin{cousinweb}
\textbf{\ml{ഇവിടെ}} (\emph{iviṭe}) is native Dravidian: the near-pointing base of \emph{itŭ} with a tail meaning \textquotedblleft{}place\textquotedblright{}. DEDR 410 gathers the near-pointing words, with Tamil \emph{iṅkē}, \textquotedblleft{}here\textquotedblright{}.
\end{cousinweb}

\subsection*{Your turn}
Say these aloud:

\begin{itemize}
\raggedright
  \item \textquotedblleft{}\ml{ഇവിടെ}\textquotedblright{} three times, picturing something different each time
  \item it once with a word you already know after it
\end{itemize}

\begin{tcolorbox}[breakable,colback=teal!4,colframe=teal!35!black,title={Before you move on}]
What is \emph{iviṭe}? (here — where I am.) What sound does it start with, and does that sound mean near, far, or a question? (\textbf{i-} — near.)
\end{tcolorbox}
\section[\texorpdfstring{\ml{അവിടെ} (aviṭe)}{അവിടെ (aviṭe)}]{\ml{അവിടെ} (aviṭe) — there — where I am not}

\label{lesson:ML-C41-there}



\begin{quote}
You can already name things. This is how you POINT at them.
\end{quote}

\subsection*{You'll want to know: \ml{അവിടെ}}
\textbf{\ml{അവിടെ}} — \emph{aviṭe} — there — where I am not.

Say it, and picture the thing it lands on as you say it. That is the whole word: it does not mean anything on its own, it means whatever the speaker is pointing at.

Its partner is \textbf{\ml{ഇവിടെ}} \emph{iviṭe}, which you have just met. The two of them differ by one sound at the front — \textbf{a-} for far — and that is not a coincidence. It is the whole system, and the last lesson of this chapter says so.

\begin{cousinweb}
\textbf{\ml{അവിടെ}} (\emph{aviṭe}) is native Dravidian: the far-pointing base of \emph{atŭ} with a tail meaning \textquotedblleft{}place\textquotedblright{}. DEDR 1 gathers the far-pointing words, with Tamil \emph{aṅkē}, \textquotedblleft{}there\textquotedblright{}.
\end{cousinweb}

\subsection*{Your turn}
\begin{itemize}
\raggedright
  \item \emph{Say it:} \textquotedblleft{}\ml{അവിടെ}\textquotedblright{} three times, picturing something different each time
  \item \emph{Say it:} it once with a word you already know after it
  \item \emph{Read it:} \textbf{\ml{ഹൃദയം}}
\end{itemize}

\begin{tcolorbox}[breakable,colback=teal!4,colframe=teal!35!black,title={Before you move on}]
What is \emph{aviṭe}? (there — where I am not.) What sound does it start with, and does that sound mean near, far, or a question? (\textbf{a-} — far.)
\end{tcolorbox}
\section[\texorpdfstring{\ml{ആര്} (ārŭ)}{ആര് (ārŭ)}]{\ml{ആര്} (ārŭ) — who? — asking about a person}

\label{lesson:ML-C41-who}



\begin{quote}
You can already name things. This is how you ASK about them.
\end{quote}

\subsection*{You'll want to know: \ml{ആര്}}
\textbf{\ml{ആര്}} — \emph{ārŭ} — who? — asking about a person.

Say it with an asking voice. It points at no one: it asks your listener to name the person.

\emph{ārŭ} opens with a long \textbf{ā-}. Do not hear it as the short \textbf{a-} of far that \textbf{\ml{അത്}} \emph{atŭ} and \textbf{\ml{അവിടെ}} \emph{aviṭe} share: it has nothing to do with distance. Tamil, next door, says \emph{yār}; Malayalam's word is the same one with its first \emph{y-} worn away.

\begin{cousinweb}
\textbf{\ml{ആര്}} (\emph{ārŭ}) is native Dravidian, the question stem \emph{yā-} worn down to \emph{ā-}. DEDR 5151 lists it with Tamil \emph{yār}, \emph{ār} and Kannada \emph{yār}, \emph{ār}.
\end{cousinweb}

\subsection*{Your turn}
\begin{itemize}
\raggedright
  \item \emph{Say it:} \textquotedblleft{}\ml{ആര്}\textquotedblright{} three times, asking about a different person each time
  \item \emph{Say it:} it once with a word you already know after it
  \item \emph{Recall:} say \emph{neñcŭ}
\end{itemize}

\begin{tcolorbox}[breakable,colback=teal!4,colframe=teal!35!black,title={Before you move on}]
What is \emph{ārŭ}? (who? — asking about a person.) Does its long ā- mean far, like the short a- of \emph{atŭ}? (\textbf{No} — it is the same word as Tamil's \emph{yār}, its y- worn away; it asks about a person.)
\end{tcolorbox}
\section[\texorpdfstring{\ml{എവിടെ} (eviṭe)}{എവിടെ (eviṭe)}]{\ml{എവിടെ} (eviṭe) — where? — asking about a place}

\label{lesson:ML-C41-where}



\begin{quote}
You can already name things. This is how you ASK about them.
\end{quote}

\subsection*{You'll want to know: \ml{എവിടെ}}
\textbf{\ml{എവിടെ}} — \emph{eviṭe} — where? — asking about a place.

Say it with an asking voice. It points at no place: it asks your listener to point one out.

Its partners are \textbf{\ml{ഇവിടെ}} \emph{iviṭe} and \textbf{\ml{അവിടെ}} \emph{aviṭe}, which you have already met. The three differ by one sound at the front — \textbf{e-} for ask — and that is not a coincidence. It is the whole system, and the last lesson of this chapter says so.

\begin{cousinweb}
\textbf{\ml{എവിടെ}} (\emph{eviṭe}) is native Dravidian: the old asking base with the same \textquotedblleft{}place\textquotedblright{} tail as \emph{iviṭe} and \emph{aviṭe}. DEDR 5151 gathers the asking words, with Tamil \emph{yāṇṭu}, \textquotedblleft{}where\textquotedblright{}.
\end{cousinweb}

\subsection*{Your turn}
\begin{itemize}
\raggedright
  \item \emph{Say it:} \textquotedblleft{}\ml{എവിടെ}\textquotedblright{} three times, asking about a different place each time
  \item \emph{Say it:} it once with a word you already know after it
  \item \emph{Read it:} \textbf{\ml{ചായ}}
  \item \emph{Recall:} say \emph{kāppi}
\end{itemize}

\begin{tcolorbox}[breakable,colback=teal!4,colframe=teal!35!black,title={Before you move on}]
What is \emph{eviṭe}? (where? — asking about a place.) What sound does it start with, and does that sound mean near, far, or a question? (\textbf{e-} — ask.)
\end{tcolorbox}
\section[i- / a- / e-]{i- / a- / e- — six words, one machine}

\label{lesson:ML-C41-deixis-system}



\begin{quote}
Say the six words from this chapter out loud, in order. Listen to how they begin.
\end{quote}

\begin{grammarlens}[title={Grammar lens: one slot changes}]
Five of these six words are \textbf{three beginnings}, reused; only the person question stands apart.

\noindent
\begin{tabularx}{\linewidth}{@{}>{\raggedright\arraybackslash}X>{\raggedright\arraybackslash}X@{}}
\toprule
\textbf{meaning} & \textbf{word} \\
\midrule
near & \textbf{\ml{ഇവിടെ}} \emph{iviṭe} \\
far & \textbf{\ml{അവിടെ}} \emph{aviṭe} \\
question & \textbf{\ml{എവിടെ}} \emph{eviṭe} \\
\bottomrule
\end{tabularx}

Malayalam spells the pattern out longest — \ml{ഇവിടെ}, \ml{അവിടെ}, \ml{എവിടെ} — which makes it the easiest to hear.

Change the front, and the meaning walks from \textbf{near} to \textbf{far} to \textbf{a question}. Nothing else in the word moves.

The same three beginnings run the thing-words too: \textbf{\ml{ഇത്}} \emph{itŭ} and \textbf{\ml{അത്}} \emph{atŭ} start exactly like the first two rows. This chapter did not teach you the matching question word for things — that one arrives later, and when it does you will already know what it has to start with.

This is worth more than the six words are. When you meet a new word in this family later, you will not have to be taught all three of it — you will be taught one and work out the others.
\end{grammarlens}

\subsection*{Your turn}
\begin{itemize}
\raggedright
  \item \emph{Say it:} the near word, then the far word, then the question word
  \item \emph{Say it:} them again for place, not for things
  \item \emph{Notice:} the part that did NOT change
\end{itemize}

\begin{tcolorbox}[breakable,colback=teal!4,colframe=teal!35!black,title={Before you move on}]
Which beginning means near? (\textbf{i-}) Far? (\textbf{a-}) And which one turns the word into a question? (\textbf{e-})
\end{tcolorbox}
